\begin{table}
\begin{tcolorbox}[tabularx={X|p{2cm}|X|},enhanced,fonttitle=\bfseries\large,fontupper=\normalsize\sffamily,
colback=yellow!10!white,colframe=red!50!black,colbacktitle=blue!30!white,
coltitle=black,title=Creating Linear Spaces of Delandtsheer-Doyen Type]
		\hline
		Command & Arguments & Purpose \\
		\hline
		\hline
\verb'-d1' & d1  &  Affine group parameter.  \\
\hline
\verb'-d2' & d2  &  Affine group parameter. \\
\hline
\verb'-q1' & q1  &  Affine group parameter. \\
\hline
\verb'-q2' & q2  &  Affine group parameter. \\
\hline
\verb'-group_label' & label  &  \\
\hline
\verb'-mask_label' & label  &  \\
\hline
\verb'-problem_label' & label  &  \\
\hline
\verb'-DDx' & x  &  Deladtsheer-Doyen parameter $x$. \\
\hline
\verb'-DDy' & y  &  Secondary Deladtsheer-Doyen parameter $y$. \\
\hline
\verb'-K' & k  & Specify the size of the base-line to generate the linear space under the affine group. \\
\hline
\verb'-pair_search_control' & label  & Options for the classification of pair orbits, see Tab.~\ref{tab:pc:options}.  \\
\hline
\verb'-search_control' &  label & Options for the classification of the base lines by means of subsets, see Tab.~\ref{tab:pc:options}. The depth of the search for starters (partial base-lines) is set here, using the \verb'-depth' command. \\
\hline
\verb'-R' & nb\_row\_types row\_type  &  \\
\hline
\verb'-C' & nb\_col\_types col\_type  &  \\
\hline
\verb'-nb_orbits_on_blocks' & n  &  \\
\hline
\verb'-masktest' & options  & Specify conditions on the mask, which is the pattern induced by the base-line when considering the grid structure. \\
\hline
\verb'-search_partial_base_lines' &   & 
Search for partial base-lines. 
The size of the partial base-line is specified in the search\_control data structure, using the \verb'-depth' argument. 
Once the classification of partial base-lines is finished (of size $s$, say), the covering matrix of pair orbits 
is written. For each orbit on partial base-lines of size $s$, 
the ${s \choose 2}$ disjoint pair-orbits orbits covered by that representative are written. \\
\hline
\verb'-singletons' & sz  & 
Search for singletons. This requires that a classification of partial base-lines of size sz has been completed before. 
This command initiates the completion of the partial base-lines. \\
\hline
\verb'-subgroup' & gens order  & Specify generators for the subgroup $G$ which acts line-transitive and point-imprimitive on the linear space. \\
\hline
\verb'-search_wrt_subgroup' &   &  \\
\hline
		\end{tcolorbox}
	\caption{\label{tab:delandtsheer:doyen}Creating Linear Spaces of Delandtsheer-Doyen Type}
	\end{table}
